\begin{figure}[tbp]
\centering
\begin{tikzpicture}[x=1.65cm]
\draw[->,viewgray] (.5,0)--(6.6,0) node[right] {$q$};
\foreach \scale in {1,3,5} {
 \node[vseed] at (\scale,0) {};
 \node[below] at (\scale,-.2) {$\scale$};
 \node[vieworange] at (\scale,.65) {$(-2,0)$};
}
\foreach \scale in {2,4,6} {
 \node[vpoint] at (\scale,0) {};
 \node[below] at (\scale,-.2) {$\scale$};
 \node[viewblue] at (\scale,.65) {$(0,0)$};
}
\end{tikzpicture}
\medskip
\begin{tabular}{@{}llll@{}}
\toprule
$x=(1/2,0)$ & $q$ odd & $q$ even & What is held fixed?\\\midrule
$\Fill_{q,x}(1)$ & $1/2$ & $+\infty$ & parameter $x$\\
$\delta_{qx}(1)$ & nonzero & zero & not the powered system\\
$\delta_x(1)$ & nonzero & nonzero & original system $\mathcal L_x$\\\bottomrule
\end{tabular}
\caption{Exact torsion behavior, not a numerical approximation. Labels
above the axis are boundary rows. Tensor powers alternate between a
nontrivial and a trivial character, while the unpowered connecting class
is unchanged.}
\label{fig:arithmetic-torsion}
\end{figure}